\section{最小假设与可复用定理模板}

\subsection{最小假设分层（审计表）}

\begin{itemize}
  \item \textbf{A0（可测动力系统）}：$(X,\mathcal{B},\mu,T)$ 为测度保持系统。
  \item \textbf{A1（有限读出）}：$\rho:X\to\{0,1\}$ 可测；窗口长度 $m$ 固定。
  \item \textbf{A2（稳定化核）}：$\Fold_m$ 可计算、满射，并实现“无相邻 $1$”的规范形（Zeckendorf 逻辑）。
  \item \textbf{A3（统计口径）}：遍历性 + 生成分割（用于调用 SMB/KS）。
\end{itemize}

其余结构（例如 cut-and-project 具体几何实现、黄金取向最优性、Hecke/\(\zeta\) 接口、HPA 对齐）均应作为额外假设/扩展端口单独列出，不得隐含进入 A0--A3。

\subsection{可选扩展假设：多尺度一致性（尺度接口）}

\begin{assumption}[A4（尺度相容/因子化）]\label{assump:A4-scale}
当需要比较不同窗口长度 $m$ 的输出时，必须额外给出尺度接口（scale interface）：对任意 $m_2\ge m_1$，存在自然投影（restriction）$\pi_{m_2\to m_1}$ 使得稳定类型空间 $\{X_m^{\mathrm Z},\pi_{m_2\to m_1}\}$ 形成逆系统，并且读出满足因子化关系
\[
\pi_{m_2\to m_1}\circ \mathcal{R}_{O,m_2}=\mathcal{R}_{O,m_1}\qquad (\mu\text{-a.e. 或逐点}).
\]
在该意义下，“结果一致”应理解为跨尺度相容性；不同 $m$ 的差异对应于可见前缀长度与由此诱导的条件结构差异。
\end{assumption}

\begin{remark}[主文固定带宽模板]
在以秩 $6$ 超空间闭合为模型类限定的论文中，主文可固定 $m=6$ 并将一般 $m$ 的可变参数化下沉到附录/协议章节。固定 $m$ 时的分辨率提升应优先通过递归读出（分割细化）表达，而不依赖提高 $m$。
\end{remark}

\subsection{模板 T1：两值/三值返回时间谱（环面诱导）}

当底系统是无理旋转且读出窗口是区间时，诱导首返结构给出有限值返回时间谱（至多三值），两值退化对应诱导为 $2$-IET，从而与旋转测度共轭。该模板在本目录用于把“读出协议的稳定条件”与“连分数/重整化更新”对齐。

\subsection{模板 T2：Sturmian 最小复杂度（环因子二值符号化）}

当 $\rho$ 来自圆周两区间分割对无理旋转的编码时，得到 Sturmian 过程，其复杂度函数满足 $p(n)=n+1$。
该模板在几何/协议论文中作为“最小复杂度读出证书”出现。

\subsection{模板 T3：Zeckendorf 合法语言计数与折叠}

合法语言 $X_m^{\mathrm Z}$ 的计数满足 Fibonacci 递推并给出 $\card{X_m^{\mathrm Z}}=F_{m+2}$；
折叠核 $\Fold_m$ 给出 $2^m\to F_{m+2}$ 的可审计压缩与纤维退化结构。
该模板同时支持：
\begin{itemize}
  \item “存在/稳定”作为投影条件化（可见态）；
  \item “不可见自由度”作为纤维对象；
  \item “边界/循环不一致”作为线性窗口语义与闭合语义的差分（Lucas/Fibonacci 恒等式）。
\end{itemize}

\subsection{模板 T4：闭包与固定点语义（现实集合）}

在 $\mathcal{P}(Y_m)$ 上定义闭包
\[
\Phi(S)=\Fold_m^{-1}(\Fold_m(S)),
\]
则固定点集合恰为纤维并。
加入可逆微观更新 $U$ 后，合成算子 $\Psi=\Phi\circ U$ 的固定点给出“闭包—演化自洽现实”的离散原型。

\subsection{模板 T5：信息论时间（SMB/KS）}

在 A3 条件下，柱集后验测度指数收缩率等于 KS 熵率，用于把“时间箭头”严格化为可计算的信息率，而非外加参数。

\subsection{模板 T6：固定带宽递归观测的验证口径}给定目标量 $G$（例如更细的局部补丁类型、转移核谱、首次到达时间分布等），在固定 $m=6$ 下构造递归输出 $x^{(0:K)}$，并用条件不确定性曲线
$$
H(G\mid x^{(0:K)})
$$
随 $K$ 的下降作为验证指标。可进一步设置对照：一次性增大 $m$ 的实证比较（仅作性能对照，不宣称理论等价）。